% ================================================================
% ch03_gyan_deepak.tex — अध्याय ३: ज्ञान दीपक
% ================================================================

\cleardoublepage
\section*{अध्याय ३: ज्ञान दीपक}
\addcontentsline{toc}{section}{अध्याय ३: ज्ञान दीपक}

\subsection*{३.१ ज्ञान के दीपक आ प्रकाश}
\addcontentsline{toc}{subsection}{३.१ ज्ञान के दीपक आ प्रकाश}

\textbf{मूल भोजपुरी रचना:}

\begin{quote}
\addfontfeatures{Color=3D3428}
ज्ञान दीपक जरत निरंतर, अज्ञान तिमिर हरत जग भर।\\
दरिया कहत सुनो मन लाई, सतगुरु बिना मुक्ति न पाई॥
\end{quote}

\textbf{भोजपुरी भावार्थ:} ज्ञान के दीपक निरंतर जरत बा आ एह संसार के अज्ञान के अन्हरिया के मिटावत बा। दरिया साहब कहत बाड़ें कि मन लगा के सुनो — सतगुरु के बिना मोक्ष ना मिल सके।

\begin{sloppypar}
{\engfont \textit{English Translation:} The Lamp of Knowledge burns continuously, dispelling the darkness of ignorance throughout the world.
Dariya says, listen with your whole mind --- liberation cannot be attained without the True Master.}
\end{sloppypar}

\begin{center}
{\color{bordercolor}\rule{0.4\textwidth}{0.4pt}}
\end{center}

\subsection*{३.२ शब्द आ सुरति अभ्यास}
\addcontentsline{toc}{subsection}{३.२ शब्द आ सुरति अभ्यास}

\textbf{मूल भोजपुरी रचना:}

\begin{quote}
\addfontfeatures{Color=3D3428}
अनहद शब्द निरंतर बाजै, सुरति निरति मिलि आतम छाजै।\\
दरिया खोजत घट ही पावा, निर्मल रूप अमर पद पावा॥
\end{quote}

\textbf{भोजपुरी भावार्थ:} भीतर अनहद नाद के शब्द निरंतर गूँजत बा। जब ध्यान ओकरा में लीन हो जाला त आत्मा आनंदित हो जाले। दरिया साहब कहत बाड़ें कि खोजला पर ई सत्य आपन भीतर ही मिल गइल।

\begin{sloppypar}
{\engfont \textit{English Translation:} The boundless inner melody resonates continuously within.
When awareness merges with this divine resonance, the soul shines in eternal peace.
Dariya found this truth within himself; the pure, immortal state was attained.}
\end{sloppypar}

\begin{center}
{\color{bordercolor}\rule{0.4\textwidth}{0.4pt}}
\end{center}

\subsection*{३.३ सतगुरु महिमा आ महत्व}
\addcontentsline{toc}{subsection}{३.३ सतगुरु महिमा आ महत्व}

\textbf{मूल भोजपुरी रचना:}

\begin{quote}
\addfontfeatures{Color=3D3428}
सतगुरु मिले त मारग पावै, अज्ञान अन्हरिया तुरंत नसावै।\\
बिनु गुरु ज्ञान न उपजै भाई, कितना कोटि उपाय कराई॥
\end{quote}

\textbf{भोजपुरी भावार्थ:} जब सतगुरु मिलेलान त सच्चा मार्ग लउके लागेला आ अज्ञानता के अन्हरिया तुरंत खतम हो जाला। बिना गुरु के कवनो ज्ञान पैदा नइखे हो सकत, चाहे कवनो करोड़ों उपाय काहे न कर लिहल जाव।

\begin{sloppypar}
{\engfont \textit{English Translation:} When the True Master is found, the path becomes clear, instantly dispelling darkness.
Without wisdom guided by Truth, spiritual awakening cannot arise, no matter how many efforts one makes.}
\end{sloppypar}

\begin{center}
{\color{bordercolor}\rule{0.4\textwidth}{0.4pt}}
\end{center}

\subsection*{३.४ विवेक आ विचार (हंस रूप पहचान)}
\addcontentsline{toc}{subsection}{३.४ विवेक आ विचार (हंस रूप पहचान)}

\textbf{मूल भोजपुरी रचना:}

\begin{quote}
\addfontfeatures{Color=3D3428}
हंस रूप सोई पहचाना, सार शब्द जिन साँचा जाना।\\
काँच कंचन में भेद जनावे, सोई सुजान परम पद पावे॥
\end{quote}

\textbf{भोजपुरी भावार्थ:} हंस के समान विवेकशील उही व्यक्ति ह जे सार शब्द आ सत्य के पहचान लेला। जे काँच आ सोना में भेद समझ लेला, उहे बुद्धिमान मोक्ष आ परम पद पावेला।

\begin{sloppypar}
{\engfont \textit{English Translation:} Like a discerning swan, true is the seeker who recognizes the essence of Truth.
One who distinguishes glass from gold attains the highest spiritual fulfillment.}
\end{sloppypar}
